%% kapa: exam=1 round=29 year=2018 period=2 form=A subject=appraisal_law subject_ko=감정평가관계법규 seq=10 number=10
%% source: Q-Net 감정평가사 기출문제 (cst00309) artlSeq=5202808 file=extracted/1008/5202808/2018년 제29회 감정평가사 1차시험 기출문제/2018년 제29회 감정평가사 1차시험 기출문제/2교시 A형.pdf
%% conversion: kordoc 4.24.51 -> tools/gen1.py (auto); review=proofread-sonnet; flags=none
\begin{question}{10}
국토의 계획 및 이용에 관한 법령상 개발행위의 허가에 관한 설명으로 옳지 않은 것은? (단, 조례는 고려하지 않음)
\begin{choices}
  \item 허가권자는 개발행위허가의 신청 내용이 도시ㆍ군계획사업의 시행에 지장이 있는 경우에는 개발행위허가를 하여서는 아니된다.
  \item 개발행위허가를 받은 부지면적을 3퍼센트 확대하는 경우에는 별도의 변경허가를 받지 않아도 된다.
  \item 성장관리방안을 수립한 지역에서 개발행위허가를 하는 경우에는 중앙도시계획위원회와 지방도시계획위원회의 심의를 거치지 아니한다.
  \item 특별시장, 광역시장, 특별자치시장, 특별자치도지사, 시장 또는 군수는 개발행위허가내용과 다르게 개발행위를 하는 자에게 그 토지의 원상회복을 명할 수 있다.
  \item 지구단위계획구역으로 지정된 지역에 대해서는 중앙도시계획위원회나 지방도시계획위원회의 심의를 거치지 아니하고 한 차례만 2년 이내의 기간 동안 개발행위허가의 제한을 연장할 수 있다.
\end{choices}
\end{question}
\answer{2}
